\tlatex
\begin{lcom}{0}%
\begin{cpar}{0}{F}{F}{0}{0}{}%
 This spec exercises operator-application \ensuremath{\.{\rightarrow}}
 TeX-macro rewriting for \ensuremath{TLATeX}.
 The target transformation from the \ensuremath{README} is, for example:
\end{cpar}%
\vshade{5.0}%
\begin{cpar}{0}{T}{F}{10.0}{0}{}%
 \ensuremath{Integral(a \.{+} b,\, c,\, d) \.{\implies}}
 \ensuremath{\backslash}int\_\ensuremath{\{c\}\.{\ct}\{d\}} (\ensuremath{a
 \.{+} b}) \ensuremath{\.{\,\backslash\,}}, dx
\end{cpar}%
\vshade{5.0}%
\begin{cpar}{1}{F}{F}{0}{0}{}%
To do this the typesetter must locate each argument by scanning forward and
 counting parentheses, commas, and the matching right-paren. The definitions
 below give every ``macro-like'' operator a well-formed body so the module
 still
 parses with \ensuremath{SANY}; only the *call sites* matter for typesetting.
\end{cpar}%
\end{lcom}%
\@x{}\moduleLeftDash\@xx{ {\MODULE} MacroOperators}\moduleRightDash\@xx{}%
\@x{ {\EXTENDS} Reals}%
\@pvspace{8.0pt}%
\@x{}%
\@y{\@s{0}%
 Operators intended to become \ensuremath{TeX} macros with sub/superscripts.
}%
\@xx{}%
\@x{ Integral ( f ,\, lo ,\, hi ) \.{\defeq} hi \.{-} lo}%
\@y{\@s{0}%
 \ensuremath{\.{\rightarrow}}
 \ensuremath{\backslash}int\_\ensuremath{\{lo\}\.{\ct}\{hi\} f
 \.{\,\backslash\,}}, dx
}%
\@xx{}%
\@x{ Sum ( term ,\, lo ,\, hi ) \.{\defeq} hi \.{-} lo}%
\@y{\@s{0}%
 \ensuremath{\.{\rightarrow}}
 \ensuremath{\backslash}sum\_\ensuremath{\{lo\}\.{\ct}\{hi\}} term
}%
\@xx{}%
\@x{ Product ( term ,\, lo ,\, hi ) \.{\defeq} hi \.{-} lo}%
\@y{\@s{0}%
 \ensuremath{\.{\rightarrow}}
 \ensuremath{\backslash}prod\_\ensuremath{\{lo\}\.{\ct}\{hi\}} term
}%
\@xx{}%
\@x{ Sqrt ( x ) \.{\defeq} x}%
\@y{\@s{0}%
 \ensuremath{\.{\rightarrow}} \ensuremath{\backslash}sqrt\ensuremath{\{x\}
}}%
\@xx{}%
\@x{ Frac ( a ,\, b ) \.{\defeq} a}%
\@y{\@s{0}%
 \ensuremath{\.{\rightarrow}}
 \ensuremath{\backslash}frac\ensuremath{\{a\}\{b\}
}}%
\@xx{}%
\@x{ Pow ( base ,\, exp ) \.{\defeq} base}%
\@y{\@s{0}%
 \ensuremath{\.{\rightarrow} base\.{\ct}\{exp\}
}}%
\@xx{}%
\@x{ Norm ( v ) \.{\defeq} v}%
\@y{\@s{0}%
 \ensuremath{\.{\rightarrow}} \ensuremath{\backslash}lVert \ensuremath{v}
 \ensuremath{\backslash}rVert
}%
\@xx{}%
\@x{ Abs ( x ) \.{\defeq} x}%
\@y{\@s{0}%
 \ensuremath{\.{\rightarrow}} \ensuremath{\backslash}lvert \ensuremath{x}
 \ensuremath{\backslash}rvert
}%
\@xx{}%
\@pvspace{8.0pt}%
\@x{}%
\@y{\@s{0}%
 Single-argument call.
}%
\@xx{}%
\@x{ A \.{\defeq} Sqrt ( 2 )}%
\@pvspace{8.0pt}%
\@x{}%
\@y{\@s{0}%
 Multi-argument calls: arguments are simple atoms.
}%
\@xx{}%
\@x{ B \.{\defeq} Integral ( 1 ,\, 0 ,\, 10 )}%
\@x{ C \.{\defeq} Frac ( 3 ,\, 4 )}%
\@x{ D \.{\defeq} Pow ( 2 ,\, 8 )}%
\@pvspace{8.0pt}%
\@x{}%
\@y{\@s{0}%
 Arguments that are themselves compound expressions containing commas-in-sets
}%
\@xx{}%
\@x{}%
\@y{\@s{0}%
 and parentheses -- the parser must not mistake the inner comma/parens for
}%
\@xx{}%
\@x{}%
\@y{\@s{0}%
 argument separators.
}%
\@xx{}%
\@x{ E \.{\defeq} Integral ( 2 \.{*} 3 \.{+} 1 ,\, 0 ,\, 10 )}%
\@x{ F \.{\defeq} Sum ( ( 1 \.{+} 2 ) \.{*} 3 ,\, 1 ,\, 100 )}%
\@x{ G \.{\defeq} Frac ( Abs ( \.{-} 5 ) ,\, Norm ( 7 ) )}%
\@pvspace{8.0pt}%
\@x{}%
\@y{\@s{0}%
 A set literal inside an argument: the commas here are NOT argument
}%
\@xx{}%
\@x{}%
\@y{\@s{0}%
 separators of \ensuremath{Product}.
}%
\@xx{}%
 \@x{ H \.{\defeq} Product ( \{ 1 ,\, 2 ,\, 3 \} \.{=} \{ 1 ,\, 2 ,\, 3 \} ,\,
 1 ,\, 5 )}%
\@pvspace{8.0pt}%
\@x{}%
\@y{\@s{0}%
 Realistic combination mirroring the \ensuremath{README} example
 \ensuremath{\backslash}int\ensuremath{\.{\ct}\{a\.{+}b\}}\_\ensuremath{\{c\}
 d}.
}%
\@xx{}%
 \@x{ Example ( a ,\, b ,\, c ,\, d ) \.{\defeq} Integral ( a \.{+} b ,\, c
 ,\, d )}%
\@pvspace{8.0pt}%
\@x{}\bottombar\@xx{}%
\@pvspace{8.0pt}%
